\documentclass{amsart}

\numberwithin{equation}{section}

%%%%%%%packages
\usepackage{amssymb}
\usepackage{enumerate, xspace}
\usepackage{dsfont}
\usepackage{showkeys}
\usepackage[bottom,first]{draftcopy}
\usepackage{changebar}
\usepackage[colorlinks]{hyperref}


%%%%%%%%%Pagination
\hfuzz=15pt

%%%%%%%%%Theorems
\newtheorem{thmm}{Theorem}
\newtheorem{thm}{Theorem}[section]
\newtheorem{lem}[thm]{Lemma}
\newtheorem{cor}[thm]{Corollary}
\newtheorem{sublem}[thm]{Sub-lemma}
\newtheorem{prop}[thm]{Proposition}
\newtheorem{conj}{Conjecture}\renewcommand{\theconj}{}
\newtheorem{defin}{Definition}
\newtheorem{cond}{Condition}
\newtheorem{hyp}{Hypotheses}
\newtheorem{exmp}{Example}
\newtheorem{rem}[thm]{Remark}

%%%%%%%%%%mathcal
\newcommand\cA{{\mathcal A}}
\newcommand\cB{{\mathcal B}}
\newcommand\cC{{\mathcal C}}
\newcommand\cD{{\mathcal D}}
\newcommand\cE{{\mathcal E}}
\newcommand\cF{{\mathcal F}}
\newcommand\cG{{\mathcal G}}
\newcommand\cH{{\mathcal H}}
\newcommand\cI{{\mathcal I}}
\newcommand\cL{{\mathcal L}}
\newcommand\cM{{\mathcal M}}
\newcommand\cN{{\mathcal N}}
\newcommand\cO{{\mathcal O}}
\newcommand\cP{{\mathcal P}}
\newcommand\cQ{{\mathcal Q}}
\newcommand\cR{{\mathcal R}}
\newcommand\cS{{\mathcal S}}
\newcommand\cT{{\mathcal T}}
\newcommand\cU{{\mathcal U}}
\newcommand\cV{{\mathcal V}}
\newcommand\cZ{{\mathcal Z}}

%%%%%%%%%%%%mathbb
\newcommand\bA{{\mathbb A}}
\newcommand\bB{{\mathbb B}}
\newcommand\bC{{\mathbb C}}
\newcommand\bD{{\mathbb D}}
\newcommand\bE{{\mathbb E}}
\newcommand\bF{{\mathbb F}}
\newcommand\bG{{\mathbb G}}
\newcommand\bH{{\mathbb H}}
\newcommand\bI{{\mathbb I}}
\newcommand\bL{{\mathbb L}}
\newcommand\bM{{\mathbb M}}
\newcommand\bN{{\mathbb N}}
\newcommand\bO{{\mathbb O}}
\newcommand\bP{{\mathbb P}}
\newcommand\bQ{{\mathbb Q}}
\newcommand\bR{{\mathbb R}}
\newcommand\bS{{\mathbb S}}
\newcommand\bT{{\mathbb T}}
\newcommand\bU{{\mathbb U}}
\newcommand\bV{{\mathbb V}}
\newcommand\bZ{{\mathbb Z}}

%%%%%%%%%%%greek
\newcommand\ve{\varepsilon}
\newcommand\eps{\epsilon}
\newcommand\vf{\varphi}
\newcommand\Var{\Sigma^2}




%%%%%%%%%%%%%%%various
\newcommand\Id{{\mathds{1}}}
\newcommand\Tr{\text{\bf Tr\,}}
\newcommand\Exp{\operatorname{Exp}}
\newcommand{\marginnote}[1]
           {\mbox{}\marginpar{\raggedright\hspace{0pt}{$\blacktriangleright$
#1}}}
\newcommand\nn{\nonumber}
\newcommand{\bigtimes}{{\sf X}}
\newcommand{\bfP}{{\bf P}}
%%%%%%%%%%%%%%%%%%Figures
\setlength{\unitlength}{1mm}

%%%%%%%%%%%%%%MikkoArvind
\newcommand{\Leb}{{\mathfrak m}}
\newcommand{\mfa}{{\mathfrak A}}

\begin{document}

\title[Quenched CLT]{Quenched CLT for random toral automorphism}

\author{Arvind Ayyer}
\address{Arvind Ayyer \\
Department of Physics \\
Rutgers University\\
136 Frelinghuysen Road\\
Piscataway, NJ 08854, USA.}
\email{ayyer@physics.rutgers.edu}


\author{Carlangelo Liverani}
\address{Carlangelo Liverani\\
Dipartimento di Matematica\\
II Universit\`{a} di Roma (Tor Vergata)\\
Via della Ricerca Scientifica, 00133 Roma, Italy.}
\email{{\tt liverani@mat.uniroma2.it}}

\author{Mikko Stenlund}
\address{Mikko Stenlund\\
Department of Mathematics \\
Rutgers University \\
110 Frelinghuysen Road \\
Piscataway, NJ 08854, USA.}
\email{mstenlun@math.rutgers.edu}



\date{\today}
\begin{abstract}
We establish a quenched Central Limit Theorem (CLT) for a smooth
observable of a random sequence of iterated hyperbolic maps.
\end{abstract}
\keywords{...}
\subjclass{...}
\maketitle




%%%PAPER
\section{introduction}
The issue of limit laws in dynamical systems has been widely explored in
the last decades and it has a clear relevance for physical applications. A
prime example of a physically relevant system is the study of the
statistical behavior of a Lorenz gas with randomly distributed obstacles.
In the case of periodic obstacles it is known that the system is
ergodic. 
\begin{changebar}
This follows from the recurrence \cite{Sim}, which in turns
follows from the CLT, \cite{Co, Sm}, which has been proved in \cite{bsc}. 
See \cite{Le} for more details and for the treatment of some (locally)
aperiodic cases.
\end{changebar}
On the contrary the random case (albeit one may na\"{\i}vely think of it as an easier
case) stands as a challenge.

If one considers the simplest case (the random position of the obstacles
is a small i.i.d. perturbation of a periodic configuration) then, by
Poincar\'e section, one is reduced to considering a random sequence of
hyperbolic symplectic maps. Yet, such a sequence of maps it is not i.i.d.
due to the presence of recollisions. Nevertheless, even disregarding the
recollision problem, also the purely i.i.d. case is poorly understood.

In this paper we address the easiest case in which such a situation
occurs:  an i.i.d. sequence of smooth symplectic maps. To make the
presentation as clear as possible we will steer away from the full
generality in which
the following results can be obtained (although we will comment on it) and
we will consider an i.i.d. sequence of linear  two dimensional toral
automorphisms.

For such a setting we will show that the time $N$ average of any smooth
zero mean observable has Gaussian fluctuations of order $\sqrt N$ for
almost every sequence of maps. Moreover, we identify the variance of such
Gaussian fluctuations.

Similar, but less complete, results are obtained in \cite{Ba} where the
Gaussian nature of the fluctuation is proven for each sequence of maps but
neither the amplitude nor the variance is investigated.

The paper is organized as following: Section \ref{sec:results} contains
the precise description of the model we will discuss and states the main
results of the paper. Such results depend on the understanding of the
ergodic properties of sequences of maps. These are investigated in Section
\ref{sec:gap} where the needed ergodic properties are related to the
spectral properties of transfer operators viewed on appropriate Banach
spaces in the spirit of the line of research started with \cite{BKL}.
Next, in Section \ref{sec:average}, we use the above results to establish
a CLT averaged over the environment for a class of systems larger than the
ones at hand but necessary to handle the quenched case. The latter is
dealt with in Section \ref{sec:quenched} using an approach inspired by the
work on random walks in random environments; see \cite{DKL} and references
therein.

\begin{changebar}
\subsection*{Acknowledgements}
We would like to thank Joel Lebowitz for posing the problem and
D. Dolgopyat for informing one of us (CL) about the Bakhtin
reference. 
\end{changebar}
MS would like to thank The Finnish Cultural Foundation for funding. MS
and AA were supported in part by NSF DMR-01-279-26 and AFOSR AF
49620-01-1-0154. 


\section{The model and the results}\label{sec:results}
Let us consider two\footnote{In fact, the following would hold almost
verbatim also for a larger collection of matrices.}  matrices
$\{A_i\}_{i=0}^{1}\in SL(2,\bN)$  and define the toral automorphisms
$T_ix=A_ix\mod 1$. Let $\{p_i\}_{i=0}^{1}\subset \bR_+$ such that
$\sum_{i=0}^1 p_i=1$. We can then introduce the Markov operator
$Q:L^\infty(\bT^2,\bR)\to L^\infty(\bT^2,\bR)$ defined by \marginnote{Let $Q$ act on $g$ and $\cL$ act on $f$ ;) ! -M}
\[
Qg(x)=\sum_{i=0}^{1}p_ig(T_i(x)).
\]
Such an operator defines a Markov Process. To describe it let us consider
the space of trajectories $\Omega_*:=(\bT^2)^{\bN}$. For each measure
$\mu$ on $\bT^2$, the above Markov process defines a probability measure
$P_\mu$ on $\Omega_*$. Let $\bE_{P_\mu}$ be the expectation with respect
to such a measure, then
\[
\begin{split}
&\bE_{P_\mu}(g(x_0))=\int g(\xi)\mu(d\xi)\\
&\bE_{P_\mu}(g(x_{i+1})\;|\; x_i)=Qg(x_i).
\end{split}
\]
Next, notice that the measure $P_\mu$ is supported on a very small set of
trajectories: if $(x_0,x_i,\dots)\in\Omega_*$ then $P_\mu$-almost surely
$x_{i+1}\in \{T_0x_i,T_{1}x_i\}$. Thus, if we consider
$\Sigma:=\{0,1\}^\bN$, we have
\[
P_\mu(\cup_{\sigma\in\Sigma}\{(x_0, T_{\sigma_1}x_0,
T_{\sigma_2}T_{\sigma_1}x_0, \dots)\})=1.
\]
In other words we can define the probability space $\Omega=\Sigma\times
\bT^2$ and the map $F:\Omega\to\Omega$ defined by
\[
F(\sigma, x)=(\tau\sigma, T_{\sigma_1}x),
\]
where $(\tau\sigma)_i=\sigma_{i+1}$, and the measure ${\bf
P}_{\wp,\mu}=\bP_{\wp}\times \mu$, where $\bP_{\wp}$ is the Bernoulli
measure with probability $\wp$ to have zero.
Finally, if we define the map $\Psi:\Omega\to\Omega_*$ by
\[
\Psi(\sigma, x):=(x, T_{\sigma_1}x, T_{\sigma_2}T_{\sigma_1}x, \dots).
\]
\begin{changebar}
It is the easy to verify that $\bE_{P_\mu}(f)=\bE_{{\bf P}_{\wp,\mu}}(f\circ
\Psi)$ for each continuos function $f$,\footnote{Inded, if $f(x_1,x_2,\dots)=g(x_n)$, then
$\bE_{P_\mu}(f)=\mu(Q^n g)=\bE_{{\bf P}_{\wp,\mu}}(f\circ \Psi)$. On
the other hand if we have already the equality for functions depending
on $n$ variables, we can write $f(x_1,\dots,x_n,x_{n+1})=g_{x_1,\dots,
x_n}(x_{n+1})$ and, by induction,
\[
\begin{split}
\bE_{P_\mu}(f)&=\bE_{P_\mu}(\bE_{P_\mu}(f\;|\; x_1,\dots, x_n))=\bE_{P_\mu}(Qg_{x_1,\dots,
x_n}(x_n))\\
&=\bE_{{\bf P}_{\wp,\mu}}\left(\sum_i p_i
f(\Psi(\sigma,x)_1, \dots, \Psi(\sigma, x)_n, T_i\Psi(\sigma,
x)_n)\right)\\
&=\bE_{{\bf P}_{\wp,\mu}}(\bE_{{\bf P}_{\wp,\mu}}(f\circ \Psi\;|\; x,
\sigma_1, \dots,\sigma_n))=\bE_{{\bf P}_{\wp,\mu}}(f).
\end{split}
\]
The assertion follows then by the density of the local functions in the
continuos one.
}
\end{changebar}
that is the two process are isomorphic hence we will use them
interchangeably as far as the study of measure theoretical properties is
concerned. Note that if $\mu$ is simultaneously $T_0$ and $T_1$ invariant,
then ${\bf P}_{\wp,\mu}$ is invariant for the map $F$. Since the maps are
symplectic, this happens for the Lebesgue measure $\Leb$. Let us set ${\bf
P}_\wp:={\bf P}_{\wp,\Leb}$.

For each function \marginnote{Why $\cC^1?$ -M Some smoothness is
needed since $f$ must be in our spaces, yet one can approximate so
exact smoothenss is not essental. -C}
$f\in\cC^1(\bT^2,\bR)$ such that $\Leb(f)=0$ we can then define the random
variables $X_k(\omega,x):=f(F^k(\omega,x))$. We are
interested in studying the $\bP_\wp$-almost sure asymptotic behavior, as
$N\to\infty$, of the random variables
\[
S_N(\omega):=\sum_{k=0}^{N-1}X_k(\omega,\cdot).
\]

The first relevant fact lies in the following Lemma.
\begin{lem}\label{lem:ergodic}
The dynamical system  $(\Omega, F, {\bf P_\wp})$ is ergodic.
\end{lem}
\begin{proof}
By the above discussion the ergodicity of $(\Omega, F, {\bf P_\wp})$ is
equivalent to the ergodicity of the stationary Markov process $P_\Leb$.
It is well known that the ergodicity of such a process is equivalent to
the fact that $Qg=g$ implies $g=\text{constant}$ for each bounded measurable $g$.
In section \ref{sec:gap} we will see (Corollary \ref{cor:trivial})
\begin{changebar}
that there exists $p,q >0$ such that,
for each $f\in\cC^{p+2d+1}, g\in\cC^{q}$, holds 
\begin{equation}\label{eq:mixing}
\lim_{n\to\infty}\Leb(fQ^ng)=\Leb(f)\Leb(g).
\end{equation}
\end{changebar}
Taking $f\in\cC^{p+2d+1}$ and $g\in  L^\infty$, we can choose a sequences
$\{g_n\}\subset\cC^{q}$ that converges to $g$ in $L^1$ and we obtain
\eqref{eq:mixing} also for such functions. But this means that $Qg=g$
implies $\Leb(f g)=\Leb(f)\Leb(g)$ for each $f\in\cC^{p+2d+1}$ which readily
implies that $g$ is constant.
\end{proof}
Since $\Leb(S_N(\omega))=0$, thanks to the previous Lemma we can apply
Birkhoff ergodic theorem and obtain
\[
\lim_{N\to\infty}\frac 1N S_N=0 \quad {\bf P_\wp}-\text{almost surely}.
\]
The next step is to investigate the variable $N^{-\frac 12}S_N$ and prove
that it satisfies a CLT. \footnote{As far as I see the following arguments
would work
verbatim in the case of ergodic markov chains but I am not sure it is
worth going to such generality. Finally the bottom line remains: given a
specific sequence it is typically impossible to decide if it is OK or
not.} Our main result is the following
\begin{thm}\label{thm:clt-quenced}
For each $\wp\in[0,1]$ there exists $\Var_\wp\in\bR_+$ such that for
$\bP_\wp$-almost all sequences $\omega$ holds\footnote{By $\cN\bigl(0,
\Var\bigr)$ we mean the centered Gaussian random variable with variance
$\Var$. The symbol $\Rightarrow$ stands for convergence in distribution.}
\[
N^{-\frac 12}S_N\Rightarrow
\cN\bigl(0, \Var_\wp\bigr)\quad under\  \Leb.
\]
\end{thm}

Before proving such a strong result we will obtain its averaged
version.
\begin{lem}\label{lem:clt-averaged}
For each $\wp\in[0,1]$ there exists $\Var_\wp\in\bR_+$ such that it holds
\[
N^{-\frac 12}S_N\Rightarrow
\cN\bigl(0, \Var_\wp\bigr)\quad under\  \bfP_\wp.
\]
\end{lem}

In turn such a result is based on a fine understanding of the dynamical
properties of certain transfer operators associated to the process. To be
more precise, for any transformation $T$ preserving the Lebesgue measure
$\Leb$ and smooth function $g$, define $\cL_{g,T}\vf:=e^{g}\,\vf\circ
T^{-1}$. In addition, given $\omega\in\Omega$ let
$\cL_{g,\omega,m}:=\cL_{g,T_{\omega_{m-1}}}\cdots\cL_{g,T_{\omega_0}}$.
Our first task will be to obtain some information on the spectral
properties of such operators, to do so in a useful way it is necessary to
introduce appropriate functional spaces. Instead of appealing to the
general theory developed in \cite{BKL, Ba1, FR, BaT, GL} we will take
advantage of the simplicity of the present setting and introduce
explicitly a particularly simple version of such theory. We will then see
how it can be used to address the ergodic theoretical questions we are
interested in.


\section{Spectral properties of the Transfer operators} \label{sec:gap}

We start to consider the case in which there is no potential, i.e. $g=1$.
Nevertheless, for further use (see section \ref{sec:quenched}) we need to
study more general automorphisms than the one introduced in the previous
section, namely $\bigoplus_{1}^d T_i\;:\;\bT^{2d}\to\bT^{2d}$ defined by
$\bigoplus_{1}^d T_ix=\bigoplus_{1}^d A_i\mod 1$ where $\bigoplus_{1}^d
A_i$
is the $d$-fold direct sums of the matrices $A_0,A_1\in SL(2,\bZ)$, that
is
\[
\bigoplus_{1}^d A_i =
\begin{pmatrix}
A_i &            & \\
   & \ddots & \\
   &            & A_i \\
\end{pmatrix}
\in SL(2d,\bZ),
\]
with $d\in\{1,2\}$,. These matrices are symplectic with respect to the
symplectic form
\[
J_{2d} =
\begin{pmatrix}
J_2 &            & \\
   & \ddots & \\
   &            & J_2 \\
\end{pmatrix},
\]
where $J_2 = \bigl( \begin{smallmatrix} 0 & 1 \\ -1 & 0\\
\end{smallmatrix} \bigr)$ is the standard symplectic form in two
dimensions.\footnote{In fact all the following can be easily generalized
to any set of symplectic toral automorphisms which preserves the standard
sector in any dimension, see \cite{LW1, LW2} for more details on the
necessary machinery.}

Let us introduce the transfer operators $\cL_{i}^{(d)} \varphi := \varphi\circ (\bigoplus_1^d T_i)^{-1}$ induced by the above automorphisms. We will consider the operator obtained from
the latter by averaging over the Bernoulli measure, namely 
\[
\cL^{(d)}_{\wp}:=\wp \,\cL^{(d)}_{0}+(1-\wp) \cL^{(d)}_{1}.
\]
Finally, we also need to study the perturbed operators given by
\[
\cL^{(d)}_{g,\wp} \varphi:=\cL^{(d)}_{\wp} (e^g \varphi).
\]

To study such operators it is necessary to introduce appropriate
Banach spaces (see \cite{GL,BaT,FR}). Here, given the simplicity of
the situation, we can use spaces inspired by
\cite{BaT,FR}.\footnote{Actually our
choice is more flexible than the one in \cite{FR}, in the spirit of
\cite{BaT}, and would allow to treat $\cC^k$
perturbations, although it is not the goal here.}

Given $v:=(v_1,\dots,v_{2d})\in\bR^{2d}$, let us denote
$\check{v}:=(v_1,v_3,\dots,v_{2d-1})\in\bR^d$ and
$\hat{v}:=(v_2,v_4,\dots,v_{2d})$. Then the standard inner product in
$\bR^d$ of the latter two reads $\langle\check{v}, \hat{v}\rangle =
\sum_{k=0}^{d-1} v_{2k+1} v_{2(k+1)}$.
Consider now the cones $\cC_-:=\{v\in\bR^{2d}\;:\; \langle\check{v},
\hat{v}\rangle \leq
0\}$, $\cC_+:=\{v\in\bR^{2d}\;:\; \langle\check{v},\hat{v}\rangle \geq
0\}$. Then
there exists $1<\lambda\leq\Lambda$ and $C_0$, depending only on
$\{A_0,A_1\}$,
such that, for all $v\in\cC_-$, $(i_1,\dots,i_n)\in \{0,1\}^n$, and
$n\in\bN$,
\begin{equation}\label{eq:contraction}
C_0^{-1}\lambda^n\leq \frac{\|(\bigoplus_1^dA_{i_n}^{-1}\cdots
\bigoplus_1^dA_{i_1}^{-1})^Tv\|}{\|v\|}\leq
C_0\Lambda^n;
\end{equation}
and, for all $v\in\cC_+$, $(i_1,\dots,i_n)\in \{0,1\}^n$, and $n\in\bN$,
\begin{equation}\label{eq:expansion}
C_0^{-1}\lambda^n\leq \frac{\|(\bigoplus_1^d A_{i_1}\cdots
\bigoplus_1^dA_{i_n})v\|}{\|v\|}\leq
C_0\Lambda^n.
\end{equation}

Moreover one can compute that there exists $\beta>1$ such that
$(\bigoplus_1^d A_i^{-1})^T\cC_-\subset \cC_\beta$, where
$\cC_\beta:=\{v\in\bR^{2d}\;:\; \beta^{-1}\|\check{v}\|^2\leq -\langle
\check{v},\hat v\rangle\leq \beta\|\check{v}\|^2\}\subset \text{int
}\cC_-$.

Now, we proceed to define the norm for the Banach space we want to
consider. Notice that the natural objects in these cones are not vectors
but Lagrangian subspaces. Recall that, given a symplectic form $J$, a
Lagrangian subspace $E\subset \bR^{2d}$ is a $d$-dimensional subspace such
that $\langle v, J w\rangle = 0$ for all $v,w \in E$.

For our choice of symplectic form, every Lagrangian subspace can also
be written as the set $E=\{ v\;:\; \hat{v} = -U \check{v}\}$ for a specific
symmetric $d\times d$ matrix $U$. Our convention is to write a minus sign in front of the $U$ here, because then $E\subset \cC_\beta$ if and only if $\beta^{-1}\Id \leq U \leq \beta \Id$.

Let us denote the set of Lagrangian subspaces as $\bL$.
For a Lagrangian subspace $E$ and a vector $k$, we set $\langle
E,k\rangle:=\displaystyle{\sup_{\substack{v\in E\\ \|v\|=1}}}\langle
v,k\rangle$.
Then, for each $p,q\in\bR_+$ and $f\in\cC^\infty(\bR^{2d},\bR)$ we
define the norm
\[
\|f\|_{p,q}:=\sup_{\substack{E \in \bL\\ E \subset
\cC_-}}\sum_{k\in\bZ^{2d}\backslash
\{0\}}|f_k|\frac{|k|^p}{1+|\langle E,k\rangle|^{p+q}}+|f_0|,
\]
where $f_k$ are the Fourier coefficients of $f$.

Notice that the $n$th power of $\cL^{(d)}_\wp$ can be expanded
\[
[\cL^{(d)}_\wp]^n = \sum_{j=1}^n\sum_{i_j=0}^1
\wp^{\delta_{i_1,0}+\dots+\delta_{i_n,0}}
(1-\wp)^{\delta_{i_1,1}+\dots+\delta_{i_n,1}} \, \cL^{(d)}_{i_1\dots i_n},
\]
where $\cL^{(d)}_{i_1\dots i_n}:= \cL^{(d)}_{i_n} \dots \cL^{(d)}_{i_1}$
.
Because the Bernoulli weights above sum to unity,
\[
\|[\cL^{(d)}_\wp]^n f\|_{p,q} \leq \sup_{(i_1,\dots,i_n)\in \{0,1\}^n}
\|\cL^{(d)}_{i_1\dots i_n} f\|_{p,q}.
\]

A straightforward computation shows that $(\cL^{(d)}_{i}
f)_k=f_{\bigoplus_1^d A_i \,k}$ for $k\in \bigoplus_1^d\bZ^2 \cong
\bZ^{2d}$. Hence, for
each $n\in\bN$ and $(i_1,\dots,i_n)\in \{0,1\}^n$,
\[
\| \cL^{(d)}_{i_1\dots i_n} f  \|_{p,q}=\sup_{\substack{E \in \bL\\ E
\subset \cC_-}}\sum_{k\in\bZ^{2d}\backslash
\{0\}}|f_k|\frac{|(\bigoplus_1^d A_{i_1}\dots \bigoplus_1^d
A_{i_n})^{-1}k |^p}{1+ |\langle E,(\bigoplus_1^d
A_{i_1}\dots \bigoplus_1^d A_{i_n})^{-1}k  \rangle
|^{p+q}}+|f_0|.
\]

We begin estimating the summand. Before that, we simplify notation and
rename the matrix product. To avoid the problem of too many indices, we
denote it only with the subscript $n$ and assume the dimension and the
sequence to be implicit:
\[
\mfa_n := \bigoplus_1^dA_{i_1}\cdots \bigoplus_1^dA_{i_n}.
\] 

Using \eqref{eq:contraction} and the fact $E \subset \cC_-$, we have
\begin{equation}\label{eq:angle0}
C_0\Lambda^n \langle  E_n,k \rangle\geq\langle E,\mfa_n^{-1} k \rangle
\geq C_0^{-1} \lambda^n \langle E_n,k \rangle \quad \text{ where } \quad
E_n:=(\mfa_n^{-1})^T E
\end{equation}
and consequently,
\[
|f_k|\frac{|\mfa_n^{-1}k |^p}{1+ |\langle E,\mfa_n^{-1}k  \rangle
|^{p+q}} \leq  C_0^{p+q} |f_k|\frac{|\mfa_n^{-1}k |^p}{1+ \lambda^{n(p+q)}
|\langle E_n,k  \rangle
|^{p+q}}.
\]

Now there are two possible cases depending on where $\mfa_n^{-1}k$ lies.
If $\mfa_{n-1}^{-1}k \notin \cC_-$, then $|\mfa_n^{-1} k| \leq C_0
\Lambda\lambda^{-n+1} |k|$ and we have
\[
\frac{|\mfa_n^{-1}k |^p}{1+ \lambda^{n(p+q)} |\langle E_n,k  \rangle
|^{p+q}} \leq C_0^p(\Lambda\lambda^{-1})^p \lambda^{-np} \frac{|k|^p}{1+
|\langle E_n,k  \rangle |^{p+q}}.
\]

While if $\mfa_{n-1}^{-1}k \in \cC_-$, then
\begin{equation}\label{eq:angle}
\langle E_n, k \rangle\geq
\frac{\lambda \Lambda^{-2n+1}|k|}{2 C_0^3 \beta\sqrt{1+\beta^2}}=:\frac{|k|}{B_{\beta,n}}.
\end{equation}
Indeed, setting $k_n:= \mfa_{n-1}^{-1}k$,
\marginnote{Had to change this back: Only $\bigl(\bigoplus_1^dA_{i_n}^{-1}\bigr)^T E\subset \cC_\beta$.}
\[
\begin{split}
\langle   \bigl(\bigoplus_1^dA_{i_n}^{-1}\bigr)^T E,  k_n \rangle&\geq 
\inf_{\substack{E\in\bL\\E\subset \cC_\beta}}\langle E, k_n \rangle
=\inf_{\beta^{-1}\Id\leq U\leq \beta\Id}\
\sup_{\check{v}\in\bR^d}\frac{\langle \check{v},\check{k}_n\rangle-\langle
\check{v},U\hat k_n\rangle}{\sqrt{|\check{v}|^2+|U\check{v}|^2}}\\
&\geq \inf_{\beta^{-1}\Id\leq U\leq \beta\Id}\frac{-\langle
\check{k}_n,\hat k_n\rangle+\langle \hat k_n,U\hat k_n\rangle}{\sqrt{|\hat
k_n|^2+|U\hat k_n|^2}}\\
&\geq \frac{|\hat k_n|}{\beta\sqrt{1+\beta^2}},
\end{split}
\]
where we have chosen $\check{v}=-\hat k_n$ in the second line. On the
other hand the choice $\check{v}=U^{-1}\check{k}_n$ yields
\[
\langle  \bigl(\bigoplus_1^dA_{i_n}^{-1}\bigr)^T E,  k_n \rangle\geq 
\inf_{\substack{E\in\bL\\E\subset \cC_\beta}}\langle E, k_n \rangle\geq 
\frac{|\check{k}_n|}{\beta\sqrt{1+\beta^2}},
\]
The inequality \eqref{eq:angle} follows from the above estimates,  \eqref{eq:angle0} and $C_0\Lambda^{n-1} |k_n|\geq |k|$.


Next, we consider two subcases. If $|k| \geq B_{\beta,n}$, then $|\langle
E_n,k \rangle |\geq 1$. Hence,
\[
\frac{|\mfa_n^{-1}k |^p}{1+ \lambda^{n(p+q)} |\langle E_n,k  \rangle
|^{p+q}} \leq C_0^p \frac{\Lambda^{pn}\lambda^{-(p+q)n} |k |^p}{\bigl
|\bigl\langle E_n,k \bigr \rangle \bigr |^{p+q}}
\leq 2C_0^p \frac{\Lambda^{pn}\lambda^{-(p+q)n} |k |^p}{1+\bigl
|\bigl\langle E_n,k \bigr \rangle \bigr |^{p+q}},
\]
which is a good estimate provided $q$ is large enough so that
$\Lambda^{p}\lambda^{-(p+q)} < 1$. The remainder is a finite sum which can
be estimated because if $E \subset \cC_-$, then $E_n \subset
(\mfa_n^{-1})^T \cC_- \subset \cC_-$. Thus,  we have
\[
\sup_{\substack{E \in \bL\\ E \subset \cC_-}}
\sum_{\substack{k\in\bZ^{2d}\backslash
\{0\}\\ |k|<B_{\beta,n}}}\frac{ |f_k|\,|\mfa_n^{-1}k |^p}{1+
\lambda^{n(p+q)} |\langle E_n,k  \rangle |^{p+q}} \leq
C_0^p \Lambda^{np} \sup_{\substack{E \in \bL\\ E \subset \cC_-}}
\sum_{\substack{k\in\bZ^{2d}\backslash
\{0\}\\ |k| < B_{\beta,n}}} \frac{|f_k|\,|k |^p}{1+ |\langle E,k  \rangle
|^{p+q}}.
\]
Accordingly, setting $\tilde \mu:=
\max\{\lambda^{-p},\Lambda^p \lambda^{-p-q}\}$ we can
collect all the above inequalities as
\begin{equation}
\label{eq:lasota-0}
\begin{split}
\|[\cL^{(d)}_\wp]^n f\|_{p,q}&\leq C_1\|f\|_{p,q},\\
\|[\cL^{(d)}_\wp]^n f\|_{p,q}&\leq 2C_0^{q+2p}\tilde \mu^n\|f\|_{p,q}+
 \sup_{\substack{E \in \bL\\ E \subset \cC_-}}
\sum_{\substack{k\in\bZ^{2d}\backslash
\{0\}\\ |k| < B_{\beta,n}}} \frac{C_0^{q+2p} \Lambda^{np} |f_k|\, |k
|^p}{1+ |\langle E,k  \rangle |^{p+q}} +|f_0|\\
&\leq C\tilde \mu^n\|f\|_{p,q}+B_n\|f\|_{p-1,q+1},
\end{split}
\end{equation}
where $B_n= C_0^{q+2p}B_{\beta,n} \Lambda^{np}$. Next, for each $\mu\in
(\tilde
\mu,1)$ choose $n_0$ such that $2C_0^{q+2p}\tilde \mu^{n_0}\leq \mu^{n_0}$
and, for
each $n\in\bN$, write $n=kn_0+m$ with $m\in\{0,\dots,n_0-1\}$. One can
thus iterate the second of the \eqref{eq:lasota-0} and obtain
\[
\|[\cL^{(d)}_\wp]^n f\|_{p,q}\leq
\mu^{kn_0}\|\cL^mf\|_{p,q}+B_{n_0}\sum_{j=0}^{k-1}\mu^{jn_0}\|\cL^{(k-1-j)n_0+m}f\|_{p-1,q+1},
\]
which finally yields
\begin{equation}
\label{eq:lasota-1}
\begin{split}
\|[\cL^{(d)}_\wp]^n f\|_{p,q}&\leq C\|f\|_{p,q},\\
\|[\cL^{(d)}_\wp]^n f\|_{p,q}&\leq C\mu^n\|f\|_{p,q}+B\|f\|_{p-1,q+1},
\end{split}
\end{equation}
with $B=CB_{n_0}(1-\mu^{n_0})^{-1}$.
We can then consider the closure, $\cB^{p,q}$, of $\cC^\infty$ in the
space of
distributions with respect to the norms $\|\cdot\|_{p,q}$. It is easy to
prove the following:


\begin{lem}
\label{lem:compactness}
The operators $\cL^{(d)}_\wp$ are well defined bounded operators on
$\cB^{p,q}$, provided $\Lambda^p<\lambda^{p+q}$. In addition, the unit
ball of $\cB^{p,q}$ is relatively compact in $\cB^{p-1,q+1}$.
\end{lem}

\begin{thm} \marginnote{Added the condition $\Lambda^p<\lambda^{p+q}$. -M}
\label{thm:quasi-compactness}
If $\Lambda^p<\lambda^{p+q}$, the operator $\cL^{(d)}_\wp$ acting on $\cB^{p,q}$ has an essential
spectral
radius smaller than $\mu$. The rest of
the spectrum consists of finitely many eigenvectors of finite
multiplicity. The only eigenvalue of modulus one is one and the constant function equal to one is the
corresponding eigenfunction.
\end{thm}
\begin{proof}
Lemma \ref{lem:compactness}, the Lasota-Yorke type inequalities
\eqref{eq:lasota-1} imply the result (see \cite{Babook, BKL} for more
details).
\end{proof}

\marginnote{The norm $\|\cdot\|_{\cC^r}$. -M} 
Before proceeding, let us mention that for any $r,n\in\bN$ the space $\cC^{r}(\bT^n,\bR)$ is endowed with the norm $\|g\|_{\cC^r} := \sum_{s=0}^r \| g^{(s)} \|_\infty$.

Moreover, a simple computation shows that $\cC^r\subset \cB^{p,q}$, provided $r>p+2d$. \marginnote{Not sure we need this. See proof of Corollary below. -M}

\begin{cor}\label{cor:trivial}
The equation \eqref{eq:mixing} holds true.
\end{cor}
\begin{proof}
\marginnote{Corrected assumptions. The proof is correct! -M}
First of all notice that, for all $f\in\cB^{p,q}$ and $g\in\cC^q$ holds
\[
\begin{split}
|\Leb(fg)|&\leq\sum_{l\in\bZ^{2}}|f_l|\,|g_{-l}|\leq \|g\|_{\cC^q} \left(|f_0|+\sum_{l\in\bZ^{2}\setminus\{0\}}|f_l|\,|l|^{-q}\right)\\
&\leq \|g\|_{\cC^q}\|f\|_{p,q}\max\left(1,\sup_{\substack{E\in\bL\\ l\in\bZ^2\setminus\{0\}}}\frac{1+|\langle E,l\rangle|^{p+q}}{|l|^{p+q}}\right)
\leq 2\|g\|_{\cC^q}\|f\|_{p,q}.
\end{split}
\]
If $p,q$ satisfy the hypothesis of Theorem \ref{thm:quasi-compactness},
$\cL^{(1)}_\wp$ has a spectral gap $\delta>0$. Thus,
\[
|\Leb(f Q^ng)-\Leb(f)\Leb(g)|=|\Leb([\cL^{(1)}_\wp]^nf \cdot
g)-\Leb(f)\Leb(g)|\leq C(1-\delta)^n\|f\|_{p,q}\|g\|_{\cC^{q}},
\]
because decomposing $\cL^{(1)}_\wp := \cQ+\cR$ with $\cQ f:=\Leb(f)$ we have $\cQ\cR=\cR\cQ=0$ and therefore $[\cL^{(1)}_\wp]^n f = \cR^n f + \cQ f$, where $\| R^n \|_{\cL(\cB^{p,q})}\leq C(1-\delta)^n$ by the spectral radius formula. \marginnote{Added this. Is it good? -M}
\end{proof}
Here is the last fact we need to know about the above functional analytic
setting.
\begin{lem}\label{lem:multiplication}
For each function $g\in\cC^{2p+q+2d+1}(\bT^{2d},\bR)$ the multiplication
operator $M_g$ defined by $M_gf=gf$, is bounded in $\cB^{p,q}$ by $C\|g\|_{\cC^{2p+q+2d+1}}$.
\end{lem}
\begin{proof}
We define the norm\footnote{We are aware that our choices of norms and the subsequent estimates, are not the optimal ones. We are simply trying to simplify the arguments as much as possible even at the expense of some, not really relevant, optimality.}
\marginnote{This proof is now correct! -M}
\[
\| g \|_{r} := \sup_{k \in \bZ^{2d}} |g_k| (1+|k|^{r}).
\]
Clearly $g\in\cC^{r}(\bT^{2d},\bR)$ implies $\|g\|_{r}\leq \|g\|_{\cC^r}$.
The $\cB^{p,q}$-norm of the product then reads
\begin{equation}\label{eq:prod-norm}
\| fg \|_{p,q} = \sup_{\substack{E \in \bL\\ E \subset \cC_-}}
\sum_{\substack{k,l\in\bZ^{2d}\\ k \neq 0}} |f_l| |g_{k-l}| \frac{ |k
|^p}{1+  |\langle E,k  \rangle |^{p+q}} + \sum_{l \in \bZ^{2d}} |f_l|
|g_{-l}|.
\end{equation}
Let us analyze the second term first:
\[
\sum_{l \in \bZ^{2d}} |f_l| |g_{-l}| \leq \| g\|_{r} \left( |f_0|
+\sum_{l \in \bZ^{2d} \setminus \{0\}} |f_l| |l|^{-r} \right).
\]
The desired bound follows, if $r\geq q$, from
\[
\frac{1}{2}| l |^{-q} \leq \frac{ | l |^p}{1+  |\langle E,l  \rangle
|^{p+q}}.
\]


Now, look at the summand in the first term of \eqref{eq:prod-norm}, ignoring
the $l=0$ case that can be taken care of separately.
\[
\begin{split}
|f_l| |g_{k-l}| \frac{ |k |^p}{1+  |\langle E,k  \rangle |^{p+q}} &\leq\|g\|_r
\left[ |f_l|  \frac{ |l |^p}{1+  |\langle E,l  \rangle |^{p+q}} \right]  \\
&\quad \times \frac{ |k |^p}{1+  |\langle E,k  \rangle |^{p+q}} \frac{1+
|\langle E,l  \rangle |^{p+q}}{ |l |^p} \frac{1}{1+|k-l|^{r}}.
\end{split}
\]
Notice that $|\langle E, l\rangle|\leq |\langle E, k\rangle|+|\langle E, k-l\rangle|$. Thus, on the one hand
\[
\sum_k \frac{ |k |^p}{1+  |\langle E,k  \rangle |^{p+q}} \frac{1+  |\langle E,k
\rangle |^{p+q}}{ |l |^p} \frac{1}{1+|k-l|^{r}} \leq \sum_k\frac{ |k |^p } {|l |^p(1+|k-l|^{r})}\leq C.
\]
One the other hand
\[
\sum_k \frac{ |k |^p}{1+  |\langle E,k  \rangle |^{p+q}} \frac{1+  |\langle E,k-l
\rangle |^{p+q}}{ |l |^p} \frac{1}{1+|k-l|^{r}} \leq 3\sum_k\frac{ |k |^p } {|l |^p(1+|k-l|^{r-p-q})}\leq C,
\]
provided $r> 2p+q+2d$. \marginnote{The majorant series has to be superharmonic ($r>\dots$ instead of $r\geq\dots$). I replaced $\cC^{2p+q+2d}$ by $\cC^{2p+q+2d+1}$ everywhere. -M} The general term of $(|\langle E, k\rangle|+|\langle E, k-l\rangle|)^{p+q}$ is bounded similarly, using $|\langle E, k\rangle|^n |\langle E, k-l\rangle|^{p+q-n} \leq (1+|\langle E, k\rangle|^n)(1+|\langle E, k-l\rangle|^{p+q-n})$ for each $n=0,\dots,p+q$.
 \end{proof}


\section{Averaged CLT}\label{sec:average}
To establish Lemma \ref{lem:clt-averaged} we must compute
\begin{changebar}
\[
\lim_{N\to\infty}\bE_{{\bf P}_\wp}\left(e^{i\frac \lambda{\sqrt
N}\sum_{k=0}^{N-1}
f_{\omega,k}}\right),
\]
where $f_{\omega,k} := X_k(\omega,\cdot) := f\circ
T_{\omega_k}\circ\dots\circ T_{\omega_0}$,
and show that this limit is the characteristic function $e^{-\frac
12\lambda^2\Var_\wp}$ of the centered normal distribution with some
variance $\Var_\wp$.
If we define the operator $\cL_{g,\wp}:=
\wp\cL_{g,T_0}+(1-\wp)\cL_{g,T_1}$, then
\begin{equation}\label{eq:powers}
\bE_{{\bf P}_\wp}\left(e^{i\frac \lambda{\sqrt N}\sum_{k=0}^{N-1}
f_{\omega,k}}\right)=
\Leb(\cL_{if\lambda N^{-\frac 12},\wp}^N\, 1).
\end{equation}
Then Theorem \ref{thm:quasi-compactness} and Lemma
\ref{lem:multiplication} show that $\cL_{if\lambda N^{-\frac 12},\wp}$ is
a bounded operator on $\cB^{p,q}$ provided $f\in\cC^{2p+q+5}$ and depends
analytically on $\lambda$. To contiune it is necessary to study the leading
eigenvalue of such an operator.
\end{changebar}

Note that, in general, given any positive operator $\cL$ on the space with maximal
simple eigenvalue one, with a spectral gap and $\Leb(\cL\vf)=\Leb(\vf)$ for
each smooth $\vf$, and any smooth complex valued function $g$ we can
define the family of operators $\cL_\nu\vf:=\cL(e^{\nu g}\vf)$ and, thanks to Lemma \ref{lem:multiplication}, the
standard perturbation theory applies. Thus there exists $\phi_\nu$,
$\mu_\nu$, with $\mu_0=1$, such that 
\[
\cL_\nu\phi_\nu=\mu_\nu \phi_\nu,\quad \Leb( \phi_\nu)=1.
\]
Differentiating this relation with respect to $\nu$ and integrating
one readily obtains
\begin{equation}\label{eq:oneder}
\mu'_\nu=\Leb \left( ge^{\nu g}\phi_\nu+e^{\nu g}\phi'_\nu \right)
\end{equation}
and, setting $\nu=0$, $\phi'_0=(\Id-\cL)^{-1}[\cL(g\phi_0)-\phi_0\Leb(
g\phi_0)]=\sum_{n=0}^\infty\cL^n[\cL(g\phi_0)-\phi_0\Leb(
g\phi_0)]$.\footnote{The latter is well defined since $(\Id-\cL)^{-1}$ is
applied on a function from which the eigendirection of $\cL$ corresponding to eigenvalue 1 has been projected out, and because of the spectral gap.}
Finally, differentiating again yields
\begin{equation}\label{eq:twoder}
\mu''_0=\Leb( g^2\phi_0)+2\sum_{n=0}^\infty\Leb( g \cL^n \left[ \cL(g\phi_0)- \phi_0\Leb(
g\phi_0)\right]). 
\end{equation}

\begin{changebar}
Thus, by standard perturbation theory and in view of Theorem \ref{thm:quasi-compactness} we can write
$\cL_\nu=\mu_\nu \cQ_\nu+\cR_\nu$ where $\cQ_\nu^2=\cQ_\nu$, $\cR_\nu\cQ_\nu=\cQ_\nu\cR_\nu=0$, the spectral radius of $\cR_\nu$ is smaller than $\rho<1$ for all $|\nu|\leq \nu_0$ for some $\nu_0\in\bR_+$ and $|\mu_\nu-1+\frac 12 \mu''_0\nu^2|\leq C|\nu|^3$ for some fixed constant $C>0$ and $|\nu|\leq \nu_0$. In addition, $\cQ_\nu$ is a rank one operator of the form $\phi_\nu\otimes m_\nu$ where $m_\nu$ belongs to the dual of the space, $m_0=\Leb$, and $|m_\nu(1)-1|\leq C|\nu|$.

If we apply the above to the operator $\cL_{if\lambda N^{-\frac 12},\wp}$, $\nu=i\lambda N^{-\frac 12}$, (hence $g=f$, and $\phi_0\equiv 1$), then remembering equation \eqref{eq:powers}
it follows that
\[
\begin{split}
\bE_{{\bf P}_\wp}\left( e^{i\frac \lambda{\sqrt
N}\sum_{k=0}^{N-1} f_{\omega,k}}\right)&=\bE_{{\bf P}_\wp}\left(\mu_\nu^N\cQ_\nu 1+\cR_\nu^N1\right)=e^{N\ln(\mu_\nu)}+\cO(|\lambda|N^{-\frac 12}+\rho^N)\\
&=e^{-\frac 12\mu''_0\lambda^2+\cO(|\lambda|^3N^{-\frac 12})}+\cO(|\lambda|N^{-\frac 12}+\rho^N).
\end{split}
\]
Hence, if $|\lambda|^3\leq \sqrt N$, we have
\begin{equation}\label{eq:clt}
\left|\bE_{{\bf P}_\wp}\left( e^{i\frac \lambda{\sqrt
N}\sum_{k=0}^{N-1} f_{\omega,k}}\right)-e^{-\frac 12\mu''_0\lambda^2}\right|\leq C \frac{|\lambda|+|\lambda|^3}{\sqrt  N}.
\end{equation}
Moreover a direct computation shows that 
\[
\Sigma_\wp^2:=\lim_{N\to\infty}\frac
1N\bE_{{\bf P}_\wp}\left(\left[\sum_{k=0}^{N-1} X_k\right]^2\right)=\mu''_0.
\]
\end{changebar}

We finish the section with two standard but important results.
\begin{lem}\label{lem:eigenvalue}
The maximal eigenvalue of $\cL^{(2)}_{i\lambda N^{-\frac 12}f_2,\wp}$ is $1-N^{-1}\lambda^2
\Var_\wp+\cO(N^{-\frac 32})$.
\end{lem}
\begin{proof}
The argument follows verbatim the previous discussion thus to prove the Lemma we only need to compute the second derivative of the leading eigenvalue. To this end we can use \eqref{eq:twoder}
which yields
\[
\mu''_0=\Leb_2( f_2^2)+2\sum_{n=0}^\infty\Leb_2\left( f_2 (\cL^{(2)})^n \cL^{(2)}(f_2)\right).
\]
which, in view of the product structure of the operator should give $2\Sigma_\wp$

\dotfill
\end{proof}

\begin{changebar}
\begin{lem}\label{lem:averaged-LDE}
Note that we can estimate, for $L>0$, \begin{equation}\label{eq:largedev}
\bfP_\wp\left(\left\{\left| \frac 1N\sum_{k=0}^{N-1}
f_{\omega,k}\right|\geq L\right\}\right)\leq Ce^{-CL^{2}N}.
\end{equation}
\end{lem}
\begin{proof}
This is an averaged large deviation estimate and can be obtained exactly
as the averaged CLT was obtained. Although the idea is standard we give here a sketch of the proof.
For any random variable $Y$ holds, for each $\beta>0$,
\[
{\bf P}_\wp(\{Y\geq L\})\leq \bE_{{\bf P}_\wp}(\Id_{\{Y\geq L\}}e^{\beta(Y-L)})\leq e^{-\beta L}\bE_{{\bf P}_\wp}(e^{\beta(Y-L)}).
\]
Applying such inequality to the present situation we have 
\[
{\bf P}_\wp(\{Y\geq L\})\leq e^{-\beta L}\Leb(\cL_{\beta (f- L)N^{-1}}^N1).
\]
Using again perturbation theory, optimizing the choice of $\beta$, and eventually applying the above argument to $-Y$, the Lemma follows.

\dotfill

\end{proof}
\end{changebar}

\section{Quenched CLT}\label{sec:quenched}
Now that we have the CLT in average we would like to establish it for a
large class of sequences. Let $\Var_\wp$ be the variance of the average
CLT
with respect to the Bernoulli process with parameter $\wp$. We wish to
show
that for $\bE_\wp$ almost all sequences $\omega$ we have the CLT with
variance $\Var_\wp$.

To this end we start with an $L^2$ estimate: assuming that $Y_N$ is a
sequence of random variables such that $\bar Y:=\lim_N \bE_\wp(Y_N)$
exists
and is real, we can compute
\[
\bE_\wp(|Y_N-\bar Y|^2)= \bE_\wp(|Y_N|^2)-\bar Y^2+2\bar Y\,\Re (\bar
Y-\bE_\wp(Y_N)).
\]
Thus, setting $f_{\omega, k}:=f\circ T_{\omega_k}\circ \cdots \circ
T_{\omega_0}$
\[
\begin{split}
\bE_\wp\left(\left|\Leb(e^{i\frac \lambda{\sqrt N}\sum_{k=0}^{N-1}
f_{\omega,k}})-e^{-\frac
12\lambda^2\Var_\wp}\right|^2\right)=&\bE_\wp\left(|\Leb(e^{i\frac
\lambda{\sqrt N}\sum_{k=0}^{N-1} f_{\omega,k}})|^2\right)\\
&\quad- e^{-\lambda^2\Var_\wp}+\cO(|\lambda|^3N^{-\frac 12}).
\end{split}
\]
The first term on the right-hand side can be conveniently reinterpreted by
introducing a product system. That is, consider the maps
$T_{\omega_k}\oplus T_{\omega_k}:\bT^4\to\bT^4$,
which are represented by the block matrices
$\Bigl(\begin{smallmatrix}A_{\omega_k} & 0 \\ 0 &
A_{\omega_k}\end{smallmatrix}\Bigr)\in SL(4,\bN)$.  Clearly they are
hyperbolic toral automorphisms (although of a higher dimensional torus)
which leave Lebesgue measure invariant. Note that all the theory of the
first section applies verbatim to this case, only the stable and unstable
directions are now two dimensional. In perfect analogy with the averaged
CLT one can define $f_2(x,y):=f(x)-f(y)$ and the operator
$
\cL^{(2)}_{i\lambda N^{-\frac 12}f_2,\wp}=\wp\cL_{i\lambda N^{-\frac 12}f_2,T_0\oplus
T_0}+(1-\wp)\cL_{i\lambda N^{-\frac 12}f_2,T_1\oplus T_1}.
$
A direct computation then shows that, calling $\Leb_2$ the Lebesgue
measure on $\bT^4$,
\[
\begin{split}
\bE_\wp\left(|\Leb(e^{i\frac \lambda{\sqrt N}\sum_{k=0}^{N-1}
f_{\omega,k}})|^2\right) & = \bE_\wp\left(\Leb_2\left (e^{i\frac
\lambda{\sqrt N}\sum_{k=0}^{N-1}
f_2\circ (T_{\omega_k}\oplus T_{\omega_k})\circ \cdots \circ
(T_{\omega_0}\oplus T_{\omega_0})} \right)\right)\\
& = \Leb_2 \bigl( \bigl [\cL^{(2)}_{i\lambda N^{-\frac 12}f_2,\wp} \bigr ]^N 1 \bigr).
\end{split}
\]
Again this problem can be treated by perturbation theory.  From
Lemma~\ref{lem:eigenvalue} follows
$\Leb_2 \bigl( \bigl [\cL^{(2)}_{i\lambda N^{-\frac 12}f_2,\wp} \bigr ]^N 1 \bigr)=(1-N^{-1}\lambda^2
\Var_\wp+\cO(|\lambda|^3N^{-\frac 
32}))^{N}= e^{-\lambda^2\Var_\wp}+\cO(|\lambda|^3N^{-\frac 12})$  \marginnote{$e^{-\lambda^2\Var_\wp}$ not $e^{-2\lambda^2\Var_\wp}$. The power is $N$, not $1/N$. -M} and therefore, remembering \eqref{eq:clt}, \marginnote{\eqref{eq:clt} lacks an error estimate for this argument. -M}
\[
\begin{split}
\bE_\wp\left(\left|\Leb(e^{i\frac \lambda{\sqrt N}\sum_{k=0}^{N-1}
f_{\omega,k}})-e^{-\frac 12
\lambda^2\Var_\wp}\right|^2\right)&= \Leb_2 \bigl( \bigl [\cL^{(2)}_{i\lambda N^{-\frac 12}f_2,\wp} \bigr ]^N 1 \bigr)-e^{-\lambda^2\Var_\wp}+\cO(|\lambda|^3N^{-\frac 12})\\
&=\cO(|\lambda|^3N^{-\frac 12}).
\end{split} 
\] 

By Chebyshev inequality the above estimate implies
\begin{equation}\label{eq:cheb}
\bP_\wp\left(\left\{\left|\Leb(e^{i\frac \lambda{\sqrt N}\sum_{k=0}^{N-1}
f_{\omega,k}})-e^{-\frac 12 \lambda^2\Var_\wp}\right|\geq
\ve\right\}\right)\leq C\ve^{-2}|\lambda|^3N^{-\frac 12}.
\end{equation}
One would then like to prove almost sure convergence by applying a
Borel-Cantelli argument but two problems are in the way: on the one hand
the sum over $N$ of the above bound diverges, on the other hand one wants
the limit to hold almost surely for all $\lambda$, that is one has
potentially uncountably many sets to deal with. Both problems can be dealt
with by applying Borel-Cantelli to subsequences and then showing that
controlling the limit of such sequences one controls the limit for each
$N$ and $\lambda$.
First of all, notice that Taylor expansion yields
\begin{equation}\label{eq:mean}
\left|e^{i\frac \lambda{\sqrt N}\sum_{k=0}^{N-1} f_{\omega,k}}-e^{i\frac
{\lambda_1}{\sqrt N}\sum_{k=0}^{N-1} f_{\omega,k}}\right|
\leq \frac {|\lambda-\lambda_1|}{\sqrt N}\left|\sum_{k=0}^{N-1}
f_{\omega,k}\right|.
\end{equation}


On the other hand, notice that the estimate \eqref{eq:largedev} in
Lemma~\ref{lem:averaged-LDE} also implies
\begin{equation}\label{eq:N}
\begin{split}
&\bfP_\wp\left(\left\{\left|\frac 1{\sqrt N}\sum_{k=0}^{N-1}
f_{\omega,k}-\frac 1{\sqrt{N+M}}\sum_{k=0}^{N+M-1} f_{\omega,k}\right|\geq
\ve\right\}\right)\\
&=\bfP_\wp\left(\left\{\left|\frac {\sqrt{1+MN^{-1}}-1}{\sqrt
{N+M}}\sum_{k=0}^{N+M-1} f_{\omega,k}-\frac
{1}{\sqrt{N}}\sum_{k=N}^{N+M-1} f_{\omega,k}\right|\geq
\ve\right\}\right)\\
& \leq \bfP_\wp\left(\left\{\left|\frac 1{N+M}\sum_{k=0}^{N+M-1}
f_{\omega,k}\right|\geq
\frac{\ve}{2\sqrt{N+M}[\sqrt{1+MN^{-1}}-1]}\right\}\right)\\
&\quad+
\bfP_\wp\left(\left\{\left|\frac 1{M}\sum_{k=0}^{M-1}
f_{\omega,k}\right|\geq \frac{\ve \sqrt{N}}{2M}\right\}\right)\leq
Ce^{-CNM^{-1}\ve^2}.
\end{split}
\end{equation}
Next, consider $b\in(\frac 12,1)$ and the sets
\footnote{Here $[x]$ stands for the integer closest to $x$.}
$A_k:=\{2^k+[j 2^{bk}]\}_{j\leq 2^{(1-b)k}}$,
$\Lambda_k:=\{-k+jk^{-1}\}_{j\leq k^2}$ and $B_k:=A_k\times \Lambda_k$.
For each $(N,\lambda)\in B_k$ let
$\Delta_k(N,\lambda)=\{(N_1,\lambda_1)\in\bN\times\bR\;:\; |N-N_1|\leq
2^{bk}+1,\, |\lambda-\lambda_1|\leq k^{-1}\}$.
Clearly
\[
\bigcup_{(N,\lambda)\in B_k}\Delta_k(N,\lambda) \supset
\left\{(N,\lambda)\in\bN\times\bR\;:\; 2^k\leq N\leq 2^{k+1},\,
|\lambda|\leq k\right\}=:J_k.
\]
We can then write
\[
\begin{split}
&\bP_\wp\left(\left\{\sup_{(N,\lambda)\in J_k}\left|\Leb(e^{i\frac
\lambda{\sqrt N}\sum_{l=0}^{N-1}f_{\omega,l}})-e^{-\frac 12
\lambda^2\Var_\wp}\right|
\geq 4\ve\right\}\right)\\
&\leq \sum_{(N,\lambda)\in
B_k}\bP_\wp\left(\left\{\sup_{(N_1,\lambda_1)\in
\Delta_k(N,\lambda)}\left|\Leb(e^{i\frac {\lambda_1}{\sqrt
N_1}\sum_{l=0}^{N_1-1}f_{\omega,l}})-e^{-\frac 12
\lambda_1^2\Var_\wp}\right|
\geq 4\ve\right\}\right)\\
&\leq\sum_{(N,\lambda)\in B_k} \Biggl\{
\bP_\wp\left(\left\{\left|\Leb(e^{i\frac {\lambda}{\sqrt
N}\sum_{l=0}^{N-1}f_{\omega,l}})-e^{-\frac 12
\lambda^2\Var_\wp}\right|\geq
\ve\right\}\right) \\
&\quad+\bP_\wp\left(\left\{\Leb\left(\sup_{(N_1,\lambda_1)\in
\Delta_k(N,\lambda)}\left|e^{i\frac {\lambda_1}{\sqrt
N_1}\sum_{l=0}^{N_1-1}f_{\omega,l}}-e^{i\frac {\lambda_1}{\sqrt
N}\sum_{l=0}^{N-1}f_{\omega,l}}\right|\right)\geq \ve\right\}\right)\\
&\quad+\bP_\wp\left(\left\{\Leb\left(\sup_{(N_1,\lambda_1)\in
\Delta_k(N,\lambda)}\left|e^{i\frac {\lambda_1}{\sqrt
N}\sum_{l=0}^{N-1}f_{\omega,l}}-e^{i\frac {\lambda}{\sqrt
N}\sum_{l=0}^{N-1}f_{\omega,l}}\right|\right)\geq \ve\right\}\right)
\Biggr\},
\end{split}
\]
where we have assumed $\Sigma_\wp k^{-1}\leq \ve$ in order to deal with
the
difference $e^{-\frac 12 \lambda^2\Var_\wp}-e^{-\frac 12
\lambda_1^2\Var_\wp}$.
 For each bounded function $g\geq 0$ holds
\[
{\bf P}_\wp(\{g \geq A\})={\bf E}_\wp(\Leb(\Id_{\{g \geq A\}}))\geq {\bf
E}_\wp(\Leb(\Id_{\{g \geq A\}}\Id_{\{\Leb(g) \geq 2A\}})).
\]
But $\Leb(g)\leq |g|_\infty\Leb(\{g\geq A\})+A$, and $\Leb (g)\geq 2A$
implies $\Leb(\{g\geq A\})\geq A|g|_\infty^{-1}$. Thus,
\[
\bP_\wp(\{\Leb(g)\geq 2A\}) \leq A^{-1} |g|_\infty{\bf P}_\wp(\{g\geq
A\}).
\]
We can estimate the above expression by
\[
\begin{split}
&\bP_\wp\left(\left\{\sup_{(N,\lambda)\in J_k}\left|\Leb(e^{i\frac
\lambda{\sqrt N}\sum_{l=0}^{N-1}f_{\omega,l}})-e^{-\frac 12
\lambda^2\Var_\wp}\right|
\geq 4\ve\right\}\right)\\
&\leq \sum_{(N,\lambda)\in
B_k}\Biggl[\bP_\wp\left(\left\{\left|\Leb(e^{i\frac {\lambda}{\sqrt
N}\sum_{l=0}^{N-1}f_{\omega,l}})-e^{-\frac 12 \lambda^2\Var_\wp}\right|
\geq \ve\right\}\right) \\
&\quad+4\ve^{-1}{\bf P}_\wp\left(\left\{\sup_{(N_1,\lambda_1)\in
\Delta_k(N,\lambda)}\left|e^{i\frac {\lambda_1}{\sqrt
N_1}\sum_{l=0}^{N_1-1}f_{\omega,l}}-e^{i\frac {\lambda_1}{\sqrt
N}\sum_{l=0}^{N-1}f_{\omega,l}}\right|\geq \frac \ve 2\right\}\right)\\
&\quad+4\ve^{-1}{\bf P}_\wp\left(\left\{\sup_{(N_1,\lambda_1)\in
\Delta_k(N,\lambda)}\left|e^{i\frac {\lambda_1}{\sqrt
N}\sum_{l=0}^{N-1}f_{\omega,l}}-e^{i\frac {\lambda}{\sqrt
N}\sum_{l=0}^{N-1}f_{\omega,l}}\right|\geq \frac \ve 2\right\}\right)
\Biggr],
\end{split}
\]
Thus, since $|e^{ix}-e^{iy}|\leq |x-y|$,
\[
\begin{split}
&\bP_\wp\left(\left\{\sup_{(N,\lambda)\in J_k}\left|\Leb(e^{i\frac
\lambda{\sqrt N}\sum_{l=0}^{N-1}f_{\omega,l}})-e^{-\frac 12
\lambda^2\Var_\wp}\right|
\geq 4\ve\right\}\right)\\
&\leq \sum_{(N,\lambda)\in
B_k}\Biggl[\bP_\wp\left(\left\{\left|\Leb(e^{i\frac {\lambda}{\sqrt
N}\sum_{l=0}^{N-1}f_{\omega,l}})-e^{-\frac 12 \lambda^2\Var_\wp}\right|
\geq \ve\right\}\right) \\
&\quad+4\ve^{-1}{\bf P}_\wp\left(\left\{\sup_{(N_1,\lambda_1)\in
\Delta_k(N,\lambda)}|\lambda_1|\,\left|\frac 1{\sqrt
N_1}\sum_{l=0}^{N_1-1}f_{\omega,l}-\frac {1}{\sqrt
N}\sum_{l=0}^{N-1}f_{\omega,l}\right|\geq \frac \ve 2\right\}\right)\\
&\quad+4\ve^{-1}{\bf P}_\wp\left(\left\{\sup_{(N_1,\lambda_1)\in
\Delta_k(N,\lambda)}\frac {|\lambda_1-\lambda|}{\sqrt
N}\left|\sum_{l=0}^{N-1}f_{\omega,l}\right|\geq \frac \ve 2\right\}\right)
\Biggr]\\
&\leq \sum_{(N,\lambda)\in
B_k}\Biggl[\bP_\wp\left(\left\{\left|\Leb(e^{i\frac {\lambda}{\sqrt
N}\sum_{l=0}^{N-1}f_{\omega,l}})-e^{-\frac 12 \lambda^2\Var_\wp}\right|
\geq \ve\right\}\right) \\
&\quad+4\ve^{-1}\sum_{|N_1-N|\leq 2^{bk}+1}{\bf
P}_\wp\left(\left\{\left|\frac
1{\sqrt N_1}\sum_{l=0}^{N_1-1}f_{\omega,l}-\frac {1}{\sqrt
N}\sum_{l=0}^{N-1}f_{\omega,l}\right|\geq \frac \ve {4k}\right\}\right)\\
&\quad+4\ve^{-1}{\bf P}_\wp\left(\left\{\frac1{\sqrt
N}\left|\sum_{l=0}^{N-1}f_{\omega,l}\right|\geq \frac {k\ve}
2\right\}\right) \Biggr]
\end{split}
\]

Then the estimates \eqref{eq:largedev}, \eqref{eq:N} and \eqref{eq:cheb}
imply, for $k\geq \Sigma_\wp\ve^{-1}$,
\[
\begin{split}
&\bP_\wp\left(\left\{\sup_{(N,\lambda)\in J_k}\left|\Leb(e^{i\frac
\lambda{\sqrt N}\sum_{l=0}^{N-1}f_{\omega,l}})-e^{-\frac 12
\lambda^2\Var_\wp}\right|
\geq 4\ve\right\}\right)\\
&\leq C\sum_{(N,\lambda)\in
B_k}\Biggl[\ve^{-2}|\lambda|^3N^{-\frac12}+\ve^{-1}\sum_{|N_1-N|\leq
2^{bk}+1}e^{-C \frac{N}{|N-N_1|}\ve^2k^{-2}}+ \ve^{-1}
e^{-Ck^2\ve^2}\Biggr]\\
&\leq Ck^2 2^{(1-b)k}\ve^{-1}\left[\ve^{-1} k^32^{-\frac
k2}+2^{bk}e^{-C2^{k(1-b)}\ve^2k^{-2}}+e^{-C\ve^2 k^2}\right],
\end{split}
\]
for which it follows that the sum over $k$ is finite. By Borel-Cantelli it
follows that the above events $\left\{\sup_{(N,\lambda)\in
J_k}\left|\Leb(e^{i\frac \lambda{\sqrt
N}\sum_{l=0}^{N-1}f_{\omega,l}})-e^{-\frac 12 \lambda^2\Var_\wp}\right|
\geq 4\ve\right\}$ happen only finitely many times with probability one.
That is, for each $\ve>0$, there exists a random variable
$N_{\ve}:\Omega\to \bN\cup \{\infty\}$, $\bP_\wp$-almost surely finite,
such
that
\[
\sup_{|\lambda|\leq \log_2N}\left|\Leb(e^{i\frac \lambda{\sqrt
{N}}\sum_{k=0}^{N-1} f_{\omega,k}})-e^{-\frac 12 \lambda^2\Var_\wp}\right|
\leq \ve \quad\text{for}\quad N\geq N_{\ve}.
\]
Here we used the fact that, for each fixed $N$, $|\lambda|\leq \log_2N$
implies $(N,\lambda)\in J_{\lfloor{\log_2 N}\rfloor}$.
Let us call $\widetilde \Omega_\ve$ the bad set of sequences, involving
$N_{\ve}=\infty$. It is an increasing set with decreasing $\ve$, such that
$\bP_\wp(\bigcup_{\ve>0} \widetilde \Omega_\ve)=\lim_{\ve\downarrow 0}
\bP_\wp(\widetilde \Omega_\ve)=0$; the bad set is independent of $\ve$.

This concludes the proof and establishes the almost sure CLT where almost
sure means that, fixing a Bernoulli measure, the set of the sequences for
which we do not have CLT has
zero measure. Note, however, that
the limit (more precisely, the variance) is not constant but depends on
$\wp$. \qed


\begin{thebibliography}{999}
\footnotesize

%%BIBLIOGRAPHY
\bibitem{Ba} Bakhtin, V.I, {\em Random processes generated by a hyperbolic
sequence of mappings. I},  Russian Acad. Sci. Izv. Math.  {\bf 44}
(1995),  no. 2, 247--279, \\
{\em Random processes generated by a hyperbolic sequence of mappings. II},
 Russian Acad. Sci. Izv. Math.  {\bf 44}  (1995),  no. 3, 617--627.
\bibitem{Babook} V.Baladi, {\em  Positive transfer operators and decay
of correlations}, Advanced Series in Nonlinear Dynamics, {\bf 16},
World Scientific (2000).
\bibitem{Ba1} V. Baladi, {\em Anisotropic Sobolev spaces and
    dynamical transfer operators: $\cC^\infty$ foliations},  Preprint.
\bibitem{BaT} V. Baladi, M. Tsujii, {\em Anisotropic H\"older
    and Sobolev spaces for hyperbolic diffeomorphisms}, to appear in Ann.
Inst. Fourier.
\bibitem{BKL} M. Blank, G. Keller, C. Liverani, {\em
    Ruelle-Perron-Frobenius spectrum for Anosov maps}, Nonlinearity, {\bf
15}:6 (2001), 1905-1973.
\bibitem{CG} J.-R. Chazottes and S. Gou\"ezel, {\em On almost-sure
versions of classical limit theorems for dynamical systems}, Probability
Theory and Related Fields 138: 195--234, 2007.
\bibitem{Co} Conze, J.-P., {\em Sur un crit�re de r�currence en dimension
2 pour les marches stationnaires, applications}. Ergodic Theory
Dynam. Systems {\bf 19} (1999), no. 5, 1233--1245. 
\bibitem{DKL} D. Dolgopyat, G.Keller, C.Liverani, {\em Random Walk in
Markovian Environment}, preprint.
\bibitem{FR}  F. Faure and N. Roy, {\em Ruelle-Pollicott resonances for
real analytic hyperbolic maps}, Nonlinearity {\bf 19} (2006), no. 6,
1233--1252.
\bibitem{GL} S. Gou\"{e}zel and C. Liverani, {\em Banach spaces adapted to
Anosov systems}, Ergodic Theory and Dynamical Systems,  26, 1,
189--217, (2006).
\bibitem{Le}
Lenci, Marco, {\em Aperiodic Lorentz gas: recurrence and ergodicity}.
Ergodic Theory Dynam. Systems {\bf 23} (2003), no. 3, 869--883. 
\bibitem{LW1} C.Liverani, M.Wojtkowski, {\em Generalization of the
Hilbert Metric to the Space of Positive Definite Matrices}, Pacific
Journal of Mathematics,  166 n. 2, pp. 339-355, (1994).
\bibitem{LW2} C.Liverani, M.Wojtkowski, {\em Ergodicity in Hamiltonian
Systems}, Dynamics Reported,  4, pp. 130-202, (1995).
\bibitem{bsc} L A Bunimovich, Ya G Sinai, N I Chernov, {\em Statistical properties of two-dimensional hyperbolic billiards}, Russ. Math. Surv., {\bf 46} (1991), no. 4, 47--106.
\bibitem{Sim} N. Simanyi. {\em Towards a Proof of Recurrence for the
Lorentz Process}. (Dynamical Systems and Ergodic Theory, Banach Center
Publications, {\bf 23}). PWN, Polish Scientific Publishers, Warsaw, 1989.
\bibitem{Sm}
Schmidt, Klaus, {\em On joint recurrence}.
C. R. Acad. Sci. Paris S�r. I Math. {\bf 327} (1998), no. 9, 837--842. 
\end{thebibliography}
\end{document}
